\section{Deployment and Research Implications}

\subsection{Throughput, availability, and recovery}
The one-worker architecture favors simple state ownership over horizontal throughput. Each batch has a bounded size and incident transitions use database row locks. This avoids two workers independently issuing duplicate requests in the tested configuration. It also makes the worker a single availability domain; during an outage, the API may accept signals while outcomes and deadlines lag until the worker recovers. A production scaling step would partition incidents by stable correlation key, retain unique constraints, and verify transaction isolation under competing workers. The local SQLite ledger uses WAL and full synchronous writes. That improves crash durability but adds write latency on each release. Batching or relaxed durability would change the accounting invariant and should be justified by measurement, not assumed safe.

The notification path is downstream of the transaction. A missed live message can be recovered by incident revision cursor. This is operationally different from the release ledger: notifications may be delivered more than once, while a release must be charged at most once. The UI should show both outcome and missing-evidence reason, with timestamp and source provenance, because an unresolved incident can be urgent. If the agent outbox is full, the gap counter records dropped initial signals; an analyst-facing completeness indicator should expose it. A full production evaluation would exercise sustained load, process kills, database failover, and long network partitions while verifying no duplicate disclosure charge and no false support.

\subsection{External validation plan}
The next empirical stage is a reproducible traffic replay with registered access to CIC-IDS2017 or an alternative dataset whose license and provenance are documented. The replay should preserve temporal order, keep each host's raw flow features local, and partition the data by defensible organizational boundaries. It should compare local-only, fixed sharing, and decision-driven policies with the same detector and predetermined threshold. Primary endpoints should include attack recall, benign false-support rate, unresolved rate, released bytes, policy units, and time from first signal to incident outcome. Repeated runs should vary seeded partition assignment, clock skew, budget, and quorum. The experiment must report the limits of synthetic partitioning even when traffic is drawn from a public dataset. A second stage with consenting organizations could then assess whether released categories are useful to analysts and whether privacy officers accept the policy.

The present attempted CIC-IDS2017 acquisition stopped when the official archive endpoint returned an HTML registration page rather than replayable flow data. No CIC-derived accuracy or throughput number appears in this manuscript. This incompleteness is a material limitation, not a minor implementation detail. It prevents claims about generalization to real traffic distributions and about the accuracy of the local detector. The current evidence is enough to show that the protocol implementation expresses a particular cost/outcome trade-off under frozen, inspectable conditions.

\subsection{Abuse cases at each trust boundary}
Table~\ref{tab:abuse} lists failure modes that a production adoption must address. The typed catalogue removes one obvious abuse path, a free-form raw-log query, but it does not make the remaining query interface harmless. An authenticated coordinator can repeatedly request permitted entries from different snapshots and combine them with outside knowledge. The rolling-hour cap slows this collection for one owner, while retention limits and audit review constrain it over longer periods. An owner might also misconfigure a policy to expose rare categories. A policy review should consider what each category reveals at the organization's actual scale, not just the integer cost assigned in the prototype.

\begin{table}[t]
\centering\caption{Illustrative abuse cases and current boundary.}
\label{tab:abuse}\scriptsize\setlength{\tabcolsep}{2pt}
\begin{tabular}{p{0.25\columnwidth}p{0.27\columnwidth}p{0.33\columnwidth}}\hline
Actor or fault & Possible effect & Current control or gap \\\hline
Repeated coordinator queries & Reconstruct activity over time & Per-snapshot tuple and rolling-hour cap; long-term inference remains\\
Compromised agent & False derived aggregate & Enrolment and audit; no log attestation\\
Duplicate network delivery & Extra debit or evidence row & Durable tuple reuse and inbox identity\\
Worker crash & Stale outcome, delayed expiry & Database state survives; one-worker availability limit\\
Forged provenance group & Artificial quorum & Enrolment mapping; collusion untested\\
Policy tightening & Old fine value already seen & Further replay denied; prior knowledge cannot be erased\\\hline
\end{tabular}
\end{table}

A malicious event producer can also target the extraction boundary. The wire schema permits only enumerated values, which limits direct prompt injection or arbitrary text from raw logs, but the aggregate computation can still be manipulated. For example, generating many requests with controlled status codes can change the failure-rate category. This is an adversarial machine-learning problem as well as a data-integrity problem. Robustness would require tests that vary attacker control over endpoints, addresses, and timing while holding the benign workload fixed. The present synthetic generator tests predetermined scenarios and does not include adaptive evasion. It should not be used to choose operational blocking actions by itself.

An insider who can alter the local SQLite ledger could remove or rewrite release rows. The transaction rules protect against ordinary process crashes and concurrent API calls, not a privileged attacker with filesystem access. A stronger audit design would append signed release receipts to an independent log and reconcile them against coordinator acknowledgements. That would introduce another retention and privacy surface: the receipt log itself records which incident prompted which entry. The trade-off should be explicit. If an organization values confidentiality more than reconstructing every prior release, local access controls and short retention may be preferable to a permanent global audit ledger.

\subsection{Concrete external-study protocol}
A follow-on study should pre-register both the case selection and decision thresholds before viewing labels. For CIC-IDS2017, the original capture days and attack windows should be preserved to avoid mixing temporally adjacent flows between training and test partitions \cite{cicdata,cicpaper}. Each simulated organization should own a disjoint set of sensors or hosts; randomly splitting individual flows from one sensor into three ``organizations'' would inflate apparent corroboration. If the dataset lacks defensible organizational boundaries, the result should be called a partitioned replay rather than a cross-organization experiment. The mapping from raw features to the five catalogue entries and the local detector threshold should be published, because both determine whether the disclosure protocol ever sees a candidate.

The replay should report a confusion matrix with a separate unresolved column. Counting unresolved as benign would hide the effect of denial and budget faults; counting all unresolved as errors would obscure the distinction between policy refusal and a wrong rule decision. A useful primary analysis is conditional: among cases that receive enough authorized evidence before the deadline, how often does the rule support or contradict correctly? A second analysis should include all cases and report the fraction left unresolved. Policy units and network bytes should be paired with both analyses. This makes it possible to see whether a lower-exposure strategy saved data by stopping after a decision or simply by failing to obtain evidence.

The external study also needs a latency model. Artificially co-located agents make round trips negligible; real partners have varying RTT, availability, and policy review delays. A controlled WAN emulation can introduce delay, loss, reordering, and partitions with recorded seeds. The evaluation should measure from the first signal's receipt to a durable terminal outcome, including expiry. It should compare a batched rich request with sequential adaptive requests under equal failure injection. A practical service-level objective may favor a somewhat richer first request if the alternative creates several long cross-organization round trips. The present local timing figures cannot answer that design choice.


\section{Local Inference and Information Boundaries}
Keeping inference local can avoid sending the authorized snapshot to an external model provider. It does not remove sensitive information from the snapshot, prevent the local process from logging it, or make a generated report safe for every recipient. Access control applies to prompts, model outputs, caches, diagnostic logs, and exported reports. The recipient scope must travel with admitted claims through the whole controlled pipeline.

A practical resource assessment should identify the model, parameter count, quantization format, runtime, hardware memory, and maximum admitted context. Prompt processing and token generation have different costs. Long evidence snapshots can consume substantial memory before the first output token, while long reports extend generation time. Concurrency increases the memory used for active contexts and may change response latency. These relationships motivate measurement; they do not justify a specific hardware recommendation without a workload.

The bounded rewrite policy gives at most two generation attempts, but input and output token limits are still necessary. A malformed or excessively large evidence bundle should be refused or summarized deterministically before generation under a documented rule. Silent truncation can remove a contradiction or a revocation notice while preserving the claim it invalidates. If context reduction occurs, the omitted evidence and resulting scope limitation need to be recorded.

The explanation service should degrade to a deterministic representation when the model is unavailable. The upstream detector and incident record remain accessible under their own permissions. A timeout must not be presented as a negative security finding. Similarly, a successful generation does not authorize remediation. Suggestions may help an analyst decide what to investigate, but executing them would require a separate workflow with its own validation and approval.

These boundaries keep the contribution focused. The paper does not propose autonomous incident response, a retrieval system with access to arbitrary logs, or a new model-serving architecture. It specifies what an explanation may assert from a bounded snapshot and how to handle failed or outdated assertions. The pilot shows a large cost and no demonstrated model benefit; broader deployments would need to reverse that finding.
